\section{Intermediate-Guided Reasoning: convertir el proceso de razonamiento en un objeto manipulable}

La segunda unidad pasa de «cuándo aparece una capacidad» a «cómo se organiza un razonamiento». El ponente aclara que Intermediate-Guided Reasoning no es un término estándar, sino una forma de agrupar una familia de métodos: hacer que el modelo genere pasos en lenguaje natural, buscar entre varios thought states, descomponer el problema de forma recursiva o entregar la representación intermedia a un intérprete de programas. La diferencia esencial es si el estado intermedio puede verificarse, revertirse y ejecutarse.

\subsection{Un umbrella term de la clase}

La diapositiva de título marca la transición; la definición propiamente dicha viene de la explicación oral: estos métodos hacen explícita la intermediate information previa a la respuesta final, de modo que una tarea compleja pueda dividirse, revisarse o delegarse a un módulo externo. Esta sección explica primero el contexto de este umbrella term en la clase y después lo divide en cuatro tipos de representación: cadena de texto, árbol de búsqueda, programa y grafo de cómputo.

\lecturefigure{slide-13-intermediate-reasoning.jpg}{Unidad de clase sobre Intermediate-Guided Reasoning}{Video de la clase de Stanford 00:07:42}

\begin{knowledgebox}{Lectura de la figura: no convertir una expresión de clase en una taxonomy fija}
Esta clase agrupa CoT, ToT, Socratic decomposition, PAL y PoT bajo un mismo paraguas para comparar su intermediate representation; la literatura académica no ha adoptado este nombre general de manera uniforme. Estas notas conservan esa organización didáctica y siguen usando el nombre original de cada artículo.
\end{knowledgebox}

\teachervoice{El ponente dice por iniciativa propia «I don't think this is actually a term». Es una teacher voice que debe conservarse: el término surge de la organización de la clase y no debe presentarse como el nombre reconocido de un campo.}

\subsection{El mecanismo básico de Chain-of-Thought}

CoT hace que el modelo genere una serie de reasoning steps en lenguaje natural antes de la respuesta final. En arithmetic, symbolic y commonsense tasks de varios pasos, permite dividir un mapeo difícil en varias generaciones condicionadas más cortas; pero generar más tokens también aumenta las oportunidades de propagar errores y de producir explicaciones que parecen razonables sin serlo.

\lecturefigure{slide-14-cot-overview.jpg}{Chain-of-Thought descompone problemas de varios pasos mediante texto intermedio}{Video de la clase de Stanford 00:08:03}

\begin{knowledgebox}{Lectura de la figura: CoT cambia el inference trace, no los pesos del modelo}
Few-shot CoT suele mostrar en el prompt ejemplos con el formato «problema→pasos→respuesta», y luego el modelo imita ese formato. Puede activar patrones algorítmicos ya presentes en el entrenamiento, pero no garantiza que cada paso corresponda a un proceso causal interno real. La mejora de la capacidad y la fidelidad de la explicación son dos criterios de aceptación distintos.
\end{knowledgebox}

Dado un problema $x$, una trayectoria intermedia $z=(z_1,\ldots,z_K)$ y una respuesta $y$, la generación con CoT puede escribirse como
\[
p_\theta(y,z\mid x)=p_\theta(z\mid x)p_\theta(y\mid x,z).
\]
$z$ representa el reasoning trace visible, $K$ el número de pasos y $y$ la respuesta final. Aumentar $p(y,z\mid x)$ no equivale a que $z$ explique fielmente el cómputo interno del modelo; solo indica que esa trayectoria textual y la respuesta tienen, en conjunto, alta probabilidad.

\subsection{Standard prompting frente a CoT prompting}

La diapositiva de ejemplos usa el problema de las pelotas de tenis para comparar la salida directa con el cálculo paso a paso. El lector no solo debe observar si la respuesta es correcta, sino también verificar si los pasos mantienen las unidades, si omiten alguna cantidad y si la respuesta final se deduce realmente de los pasos anteriores. Esta sección parte de dos prompts para el mismo problema y muestra cómo la trayectoria visible cambia la capacidad de diagnóstico, sin constituir por sí misma una proof.

\lecturefigure{slide-15-cot-examples.jpg}{Comparación entre Standard prompting y Chain-of-Thought prompting}{Video de la clase de Stanford 00:09:18}

\begin{knowledgebox}{Lectura de la figura: los pasos ofrecen una superficie de depuración}
Cuando una direct answer es incorrecta, es difícil saber en qué punto se desvió el modelo; CoT al menos expone fallas de arithmetic, symbol mapping, missing step o semantic misunderstanding. La visibilidad mejora el diagnosis, pero no mejora automáticamente la truthfulness, porque el modelo también puede obtener primero la respuesta y luego redactar la justificación.
\end{knowledgebox}

\subsection{El análisis de errores es más útil que la exactitud promedio}

La diapositiva Improving CoT divide los errores en categorías como calculator error, symbol mapping, missing step e incoherent reasoning. El ponente subraya que, aunque los modelos grandes obtienen un performance gain notable, siguen existiendo non-negligible errors; por eso se necesitan herramientas, verifiers y un mejor entrenamiento.

\lecturefigure{slide-16-cot-improvements.jpg}{Fuentes de error de Chain-of-Thought y direcciones de mejora}{Video de la clase de Stanford 00:09:42}

\begin{knowledgebox}{Lectura de la figura: cada categoría de error corresponde a una capa de corrección distinta}
El calculator error puede delegarse a un interpreter; el symbol mapping requiere una vinculación de variables más estable; el missing step puede abordarse con decomposition o search; el incoherent CoT requiere verifier, consistency o supervisión del proceso. Atribuir todos los errores a que «el modelo no es lo bastante grande» hace perder correcciones de bajo costo.
\end{knowledgebox}

\begin{warningbox}{Un reasoning trace no es una proof}
Un CoT fluido puede contener pasos irrelevantes, saltos implícitos o racionalizaciones a posteriori. En tareas de alto riesgo hay que verificar executable constraints, citar evidencia o resultados del entorno, en lugar de sustituir la corrección por la longitud del texto en lenguaje natural.
\end{warningbox}

\subsection{Cómo hacer que los modelos pequeños también usen CoT}

En clase se señala que el CoT temprano era más eficaz a partir de escalas del orden de diez mil millones de parámetros, y que los modelos pequeños suelen cometer errores de missing step y de semantic understanding. Las líneas de investigación incluyen hacer fine-tuning con traces generados por modelos grandes, distillation, reasoning formats más concisos y verifiers especializados.

\lecturefigure{slide-17-cot-smaller-models.jpg}{Líneas de investigación de Chain-of-Thought orientadas a modelos pequeños}{Video de la clase de Stanford 00:10:18}

\begin{knowledgebox}{Lectura de la figura: el problema de los modelos pequeños puede ser de capacidad de representación o de formato de supervisión}
Si el trace es demasiado largo o contiene ruido, el modelo pequeño aprende primero la plantilla superficial; si la tarea requiere un latent state que el modelo no puede representar, la destilación por sí sola tampoco basta. Conviene evaluar por separado final answer, step validity, trace length y compute, y no tratar el «trace de un modelo grande» como el material de enseñanza óptimo por naturaleza.
\end{knowledgebox}

\subsection{De dominios específicos a un reasoning format más general}

La diapositiva de Generalization cuestiona el formato fijo de CoT y su dependencia del dominio: destaca en problemas matemáticos de varios pasos, pero no todas las tareas se prestan a pasos lineales en lenguaje natural. Un método más general debe admitir distintas state representations, branching, external calls y stopping rules.

\lecturefigure{slide-18-cot-generalization.jpg}{Rigidez de formato y generalización entre dominios en Chain-of-Thought}{Video de la clase de Stanford 00:11:03}

\begin{knowledgebox}{Lectura de la figura: la representación del razonamiento debe corresponder a la estructura de la tarea}
El texto lineal es adecuado para la descomposición secuencial; las tareas de planificación pueden requerir árboles o grafos; la aritmética exacta se presta a programas; las tareas de percepción pueden requerir un state visual. La Generalization no consiste en obligar a todas las tareas a producir el mismo tipo de explicación larga, sino en elegir una intermediate representation que pueda verificarse y ejecutarse.
\end{knowledgebox}

\subsection{Tree of Thoughts: de una sola trayectoria a la búsqueda en el espacio de estados}

Tree of Thoughts hace que el modelo genere varios thoughts candidatos, evalúe el siguiente paso y, cuando es necesario, haga look ahead o backtrack. Convierte el razonamiento de un único autoregressive sample en un search procedure; con ello gana flexibility, a cambio de más inference calls y de evaluator dependence.

\lecturefigure{slide-19-tree-of-thoughts.jpg}{Evaluación de ramas y retroceso en Tree of Thoughts}{Video de la clase de Stanford 00:12:03}

\begin{knowledgebox}{Lectura de la figura: los nodos son estados intermedios y las aristas, acciones de razonamiento}
Un mismo input produce varios candidate states; el evaluator los califica y se conservan algunas ramas; las ramas fallidas pueden revertirse. Lo más importante de la figura no es su aspecto de árbol, sino que los tres papeles de generation, evaluation y selection están separados, lo que permite controlar la búsqueda de forma explícita.
\end{knowledgebox}

Si en cada nivel se expanden $b$ candidatos y la profundidad es $d$, una búsqueda completa visita como máximo
\[
1+b+b^2+\cdots+b^d=\frac{b^{d+1}-1}{b-1}
\]
nodos. $b$ es el branching factor y $d$ la reasoning depth. ToT debe controlar el costo exponencial con beam, heuristic o pruning; los errores del evaluator también pueden podar de antemano ramas correctas.

\subsection{Socratic Questioning: descomposición recursiva y recuperación de respuestas}

Socratic Questioning usa un modelo grande para plantear subproblems, los resuelve de forma recursiva y luego combina los resultados hacia arriba. Al igual que ToT, admite backtracking, pero su enfoque didáctico se acerca más a divide-and-conquer: primero se confirma que cada subproblema tenga solución y después se integran las respuestas locales en la solución del problema original.

\lecturefigure{slide-20-socratic-questioning.jpg}{Descomposición recursiva en subproblemas de Socratic Questioning}{Video de la clase de Stanford 00:12:33}

\begin{knowledgebox}{Lectura de la figura: la decomposition quality determina el límite superior}
Si los subproblemas son incompletos, se traslapan entre sí o pierden restricciones, no se obtiene la respuesta global aunque cada respuesta local sea correcta. El sistema debe verificar coverage, dependency y recomposition, en lugar de revisar solo si cada leaf parece razonable.
\end{knowledgebox}

\begin{lstlisting}[caption={Flujo de razonamiento estructurado con verificación y retroceso}]
def solve(problem, depth):
    if depth == 0 or is_atomic(problem):
        return answer_and_verify(problem)
    subproblems = propose_decomposition(problem)
    partial = [solve(item, depth - 1) for item in subproblems]
    candidate = compose(problem, partial)
    if verify(candidate, problem):
        return candidate
    return revise_decomposition(problem, partial)
\end{lstlisting}

\subsection{PAL: el modelo de lenguaje planifica, el intérprete calcula}

Program-Aided Language Models hace que el modelo genere un programa y delega la arithmetic exacta y la symbolic execution a un interpreter como Python. El LLM se encarga de convertir el lenguaje natural en una estructura ejecutable y la computadora de ejecutarla de forma determinista; esto es más confiable que pedir al modelo que simule toda la aritmética dentro del token space.

\lecturefigure{slide-21-pal.jpg}{PAL usa un programa como representación intermedia e invoca un intérprete}{Video de la clase de Stanford 00:13:09}

\begin{knowledgebox}{Lectura de la figura: división del trabajo entre razonamiento neuronal y ejecución simbólica}
El prompt muestra el problema, el program y el answer; el modelo produce code y el interpreter lo ejecuta para obtener el resultado. La ventaja es que la arithmetic es verificable y el trace es ejecutable; el riesgo es que la code generation malinterprete la semántica, invoque API peligrosas o produzca efectos secundarios fuera del sandbox.
\end{knowledgebox}

\subsection{Program of Thoughts: un scratchpad programático más general}

Program of Thoughts también hace que el modelo genere un program, pero pone énfasis en intercalar la computation con el language reasoning. En problemas matemáticos complejos, el programa conserva variables, bucles y relaciones entre funciones, lo que evita la deriva aritmética de las cadenas largas en lenguaje natural. Esta sección retoma la división del trabajo con el interpreter de PAL y se centra en cómo un scratchpad programático conserva el estado a largo plazo y la estructura composicional.

\lecturefigure{slide-22-program-of-thoughts.jpg}{Flujo de generación y ejecución de programas en Program of Thoughts}{Video de la clase de Stanford 00:13:51}

\begin{knowledgebox}{Lectura de la figura: que el programa sea correcto no significa que el problema se haya entendido correctamente}
El Interpreter solo garantiza que el programa dado se ejecute conforme a su semántica; si el modelo traduce el problema a una fórmula equivocada, el programa producirá de forma consistente una respuesta incorrecta. Por eso el pipeline necesita, como mínimo, syntax check, unit test, constraint check y result sanity check.
\end{knowledgebox}

El program-guided reasoning puede escribirse como
\[
P\sim p_\theta(P\mid x),\qquad y=\mathcal{I}(P),
\]
donde $P$ es el programa generado por el modelo, $\mathcal{I}$ es un intérprete controlado y $y$ es el resultado de la ejecución. La confiabilidad depende de dos partes, el semantic parsing $p_\theta$ y el execution sandbox $\mathcal{I}$, y no solo de la final answer accuracy.

\subsection{Faith and Fate: expresar la compositionality como un grafo de cómputo}

Al final de la clase se presenta el trabajo sobre computation graph, con énfasis en que una representación intermedia estructurada puede dividir un problema complejo en nodos verificables. El graph de la figura no es necesariamente superior a todos los métodos textuales, pero ofrece dependency, reuse y local verification, y es adecuado para tareas composicionales.

\lecturefigure{slide-23-faith-and-fate.jpg}{El compositional computation graph del trabajo Faith and Fate}{Video de la clase de Stanford 00:14:33}

\begin{knowledgebox}{Lectura de la figura: la estructura de grafo hace explícitas las dependencias}
Los nodos representan subcálculos intermedios y las aristas, dependencias; si la respuesta final es incorrecta, se puede recorrer el graph para localizar el nodo erróneo. En comparación con el CoT lineal, el graph facilita reutilizar subresultados compartidos y se presta mejor a la ejecución en paralelo; el costo es que la graph construction y la node semantics requieren, a su vez, aprendizaje o reglas.
\end{knowledgebox}

\subsection{Resumen de la sección}

Lo que comparten los Intermediate-guided methods es que hacen el estado intermedio visible, explorable o ejecutable. CoT aporta texto lineal, ToT/Socratic aportan search/decomposition, PAL/PoT aportan la ejecución de programas y el computation graph aporta una estructura de dependencias; la confiabilidad depende del verifier, del interpreter y de la retroalimentación, no de la longitud de la trayectoria.
